\hypertarget{implementation}{}
\hypertarget{complete-reproducible-implementation-architecture}{%
\chapter{Complete reproducible implementation
architecture}\label{complete-reproducible-implementation-architecture}}

\hypertarget{suggested-repository-structure}{%
\section{Suggested repository
structure}\label{suggested-repository-structure}}

\begin{verbatim}
project/
|-- analysis/
|   |-- data.py                    # parsing, joins, integrity
|   |-- factor_models.py           # OLS and frozen-beta evaluation
|   |-- joint_hac.py               # full alpha covariance
|   |-- empirical_bayes.py         # profile ML and posterior
|   |-- multiple_testing.py        # BH families
|   |-- bootstrap.py               # common-date circular blocks
|   |-- persistence.py             # overlap and transitions
|   |-- forward_validation.py      # strict forward tests
|   |-- incremental_validation.py  # rolling vs benchmark rankings
|   |-- diagnostics.py             # JB, LB, ARCH, SFA
|   |-- external_validation.py     # source and statsmodels checks
|   `-- run_all.py                 # canonical orchestration
|-- config/
|   `-- analysis.yaml
|-- data/
|   |-- raw/                       # immutable
|   `-- SOURCES.md
|-- results/
|   |-- tables/
|   |-- figures/
|   `-- run_manifest.json
|-- reports/
|   `-- generate_report.py
|-- tests/
|-- requirements-dev-lock.txt
`-- README.md
\end{verbatim}

\hypertarget{canonical-pipeline}{%
\section{Canonical pipeline}\label{canonical-pipeline}}

The pipeline runs as a linear sequence of stages, each implemented in a
module of the same name and documented in the chapter shown:
\begin{enumerate}
\tightlist
\item
  \textbf{Load and validate} the immutable panel (Chapter~4).
\item
  \textbf{Estimate and infer}: FF3 alphas (Chapter~5), the joint HAC
  covariance (Chapter~6), the joint tests and false-discovery control
  (Chapter~7), and the empirical-Bayes posteriors (Chapter~8).
\item
  \textbf{Resample and validate}: the common-date block bootstrap
  (Chapter~9), persistence and transitions (Chapter~10), and the forward
  and incremental tests (Chapters~11--12).
\item
  \textbf{Diagnose and reconcile}: residual diagnostics and FF5 sensitivity
  (Chapter~13), then external validation (Chapter~15).
\item
  \textbf{Export} the canonical tables and \textbf{generate} the report,
  every number of which is read back from those tables rather than
  recomputed.
\end{enumerate}
The whole sequence is deterministic --- fixed seed 2026 and pinned
dependencies --- so a second run reproduces the outputs bit for bit. It is
invoked by:
\begin{verbatim}
python -m pip install -r requirements-dev-lock.txt
pytest -q
ruff check .
MPLBACKEND=Agg python -m analysis.run_all
python -m analysis.export_tables
python reports/generate_report.py
\end{verbatim}

\hypertarget{run-manifest}{%
\section{Run manifest}\label{run-manifest}}

A run manifest should record:

\begin{itemize}
\tightlist
\item
  input paths, byte sizes and hashes;
\item
  sample dates, dimensions and weighting block;
\item
  model, factor columns, return units and annualisation;
\item
  HAC lag/kernel/correction and covariance diagnostics;
\item
  bootstrap type, block length, draws and seed;
\item
  rolling window, step, forward horizon and comparison benchmarks;
\item
  software/Python versions, Git commit if available, runtime and output
  hashes.
\end{itemize}

\hypertarget{essential-tests}{%
\section{Essential tests}\label{essential-tests}}

\hypertarget{data-tests}{%
\subsection{Data tests}\label{data-tests}}

\begin{itemize}
\tightlist
\item
  Duplicate key rejection
\item
  Missing month rejection
\item
  Panel balance
\item
  Factor consistency
\item
  Unit and block detection
\end{itemize}

\hypertarget{estimator-tests}{%
\subsection{Estimator tests}\label{estimator-tests}}

\begin{itemize}
\tightlist
\item
  OLS vs statsmodels
\item
  Joint HAC diagonal special case
\item
  PSD safeguards
\item
  EB shrinkage simulations
\item
  GRS/Wald null and alternative
\end{itemize}

\hypertarget{validation-tests}{%
\subsection{Validation tests}\label{validation-tests}}

\begin{itemize}
\tightlist
\item
  No-look-ahead dates
\item
  Frozen-beta invariants
\item
  Common-date bootstrap
\item
  BH edge cases
\item
  Increment sign conventions
\end{itemize}

\hypertarget{reporting-tests}{%
\subsection{Reporting tests}\label{reporting-tests}}

\begin{itemize}
\tightlist
\item
  Headline/table consistency
\item
  No p-value printed as zero
\item
  Status semantics
\item
  All claims generated from canonical tables
\item
  Readable rendered pages
\end{itemize}

\hypertarget{two-notebooks-two-purposes}{%
\section{Two notebooks, two purposes}\label{two-notebooks-two-purposes}}

A beginner benefits from separating conceptual verification from full
empirical reproduction. The teaching notebook should be tiny enough to
inspect every matrix. The replication notebook should call production
functions rather than contain a second, drifting implementation.

\begin{longtable}[]{@{}p{0.20\textwidth}p{0.34\textwidth}p{0.36\textwidth}@{}}
\toprule\noalign{}
Notebook & Suggested cells & Required outputs \\
\midrule\noalign{}
\endhead
\bottomrule\noalign{}
\endlastfoot
\textbf{01\_mathematical\_foundations.ipynb} & Simulate factors and
three portfolios; derive OLS by hand; compare matrix OLS; build
\ensuremath{\Gamma }\textsubscript{0} and \ensuremath{\Gamma }\textsubscript{1}; complete the EB square;
generate circular blocks; simulate half-normal SFA; compare
analytic/numerical gradients & Small matrices printed in full,
hand-versus-code equality assertions, conditional-density integral equal
to one, SFA unit-invariance assertion \\
\textbf{02\_full\_replication.ipynb} & Load manifest; call canonical
data validation; fit FF3; joint HAC/Wald/GRS/BH; EB; bootstrap;
persistence; forward/incremental tests; diagnostics/SFA; FF5 and source
checks & Canonical tables loaded from disk, headline reconciliation
table, hashes and run metadata; no duplicated estimator code \\
\end{longtable}

\hypertarget{a-literate-cell-pattern}{%
\section{A literate cell pattern}\label{a-literate-cell-pattern}}

\begin{verbatim}
# Question
# Does the full joint HAC diagonal reproduce the 25 separate HAC variances?

# Mathematical object
# V_HAC = c * [Gamma_0 + sum_l w_l(Gamma_l + Gamma_l.T)]

# Computation
V_joint = joint_newey_west_alpha_covariance(resid, X, lags=12)
v_separate = np.array([single_alpha_hac_variance(resid[:, i], X, 12)
                       for i in range(N)])

# Falsifiable check
np.testing.assert_allclose(np.diag(V_joint), v_separate, rtol=1e-10)

# Interpretation
# Equality validates implementation consistency; it does not validate the model.
\end{verbatim}

Each notebook section should state the question, mathematical object,
computation, falsifiable check and interpretation. This makes the
notebook an audit trail rather than a scrapbook of commands.

